% orig_front.tex -- title and abstract of the Checkpoint 2 report (BTP_Checkpoint2.pdf, p. 1),
% re-typeset verbatim. The original title block has NO author, affiliation or date lines:
% it was produced by \maketitle with empty \author{} and \date{}.
% Usage (as in the original):
%   \title{\origtitle}\author{}\date{}\maketitle
%   \begin{abstract}\origabstract\end{abstract}
% The section references in the abstract are \ref's to labels defined in orig.tex
% (they were \ref's in the original too, and they are correct there).
\newcommand{\origtitle}{A Hierarchical Task Network Formulation of Orb-Web Construction Behaviour in \textit{Araneus diadematus}}
\newcommand{\origabstract}{\noindent We formalise the orb-web construction behaviour described by Zschokke (1996) as a Hierarchical Task Network (HTN) planning problem. The model distinguishes two behavioural regimes documented empirically: an early, \textit{opportunistic} phase (exploration, bridging, and proto-hub establishment) in which the spider selects among several alternative motor strategies, and a later, \textit{stereotyped} phase (frame, radii, auxiliary spiral, capture spiral) that follows a fixed procedural sequence. Section~\ref{sec:orig-htn} states the generic forward-search HTN planning procedure; Section~\ref{sec:orig-domain} gives the domain-specific task decomposition, including the backtracking search over proto-hub construction strategies that reproduces the variability Zschokke observed; Section~\ref{sec:orig-energy} defines the locomotory energy metric used to evaluate resulting plans.}
